% Options for packages loaded elsewhere
\PassOptionsToPackage{unicode}{hyperref}
\PassOptionsToPackage{hyphens}{url}
\documentclass[
]{article}
\usepackage{xcolor}
\usepackage[margin=1in]{geometry}
\usepackage{amsmath,amssymb}
\setcounter{secnumdepth}{-\maxdimen} % remove section numbering
\usepackage{iftex}
\ifPDFTeX
  \usepackage[T1]{fontenc}
  \usepackage[utf8]{inputenc}
  \usepackage{textcomp} % provide euro and other symbols
\else % if luatex or xetex
  \usepackage{unicode-math} % this also loads fontspec
  \defaultfontfeatures{Scale=MatchLowercase}
  \defaultfontfeatures[\rmfamily]{Ligatures=TeX,Scale=1}
\fi
\usepackage{lmodern}
\ifPDFTeX\else
  % xetex/luatex font selection
\fi
% Use upquote if available, for straight quotes in verbatim environments
\IfFileExists{upquote.sty}{\usepackage{upquote}}{}
\IfFileExists{microtype.sty}{% use microtype if available
  \usepackage[]{microtype}
  \UseMicrotypeSet[protrusion]{basicmath} % disable protrusion for tt fonts
}{}
\makeatletter
\@ifundefined{KOMAClassName}{% if non-KOMA class
  \IfFileExists{parskip.sty}{%
    \usepackage{parskip}
  }{% else
    \setlength{\parindent}{0pt}
    \setlength{\parskip}{6pt plus 2pt minus 1pt}}
}{% if KOMA class
  \KOMAoptions{parskip=half}}
\makeatother
\usepackage{color}
\usepackage{fancyvrb}
\newcommand{\VerbBar}{|}
\newcommand{\VERB}{\Verb[commandchars=\\\{\}]}
\DefineVerbatimEnvironment{Highlighting}{Verbatim}{commandchars=\\\{\}}
% Add ',fontsize=\small' for more characters per line
\usepackage{framed}
\definecolor{shadecolor}{RGB}{248,248,248}
\newenvironment{Shaded}{\begin{snugshade}}{\end{snugshade}}
\newcommand{\AlertTok}[1]{\textcolor[rgb]{0.94,0.16,0.16}{#1}}
\newcommand{\AnnotationTok}[1]{\textcolor[rgb]{0.56,0.35,0.01}{\textbf{\textit{#1}}}}
\newcommand{\AttributeTok}[1]{\textcolor[rgb]{0.13,0.29,0.53}{#1}}
\newcommand{\BaseNTok}[1]{\textcolor[rgb]{0.00,0.00,0.81}{#1}}
\newcommand{\BuiltInTok}[1]{#1}
\newcommand{\CharTok}[1]{\textcolor[rgb]{0.31,0.60,0.02}{#1}}
\newcommand{\CommentTok}[1]{\textcolor[rgb]{0.56,0.35,0.01}{\textit{#1}}}
\newcommand{\CommentVarTok}[1]{\textcolor[rgb]{0.56,0.35,0.01}{\textbf{\textit{#1}}}}
\newcommand{\ConstantTok}[1]{\textcolor[rgb]{0.56,0.35,0.01}{#1}}
\newcommand{\ControlFlowTok}[1]{\textcolor[rgb]{0.13,0.29,0.53}{\textbf{#1}}}
\newcommand{\DataTypeTok}[1]{\textcolor[rgb]{0.13,0.29,0.53}{#1}}
\newcommand{\DecValTok}[1]{\textcolor[rgb]{0.00,0.00,0.81}{#1}}
\newcommand{\DocumentationTok}[1]{\textcolor[rgb]{0.56,0.35,0.01}{\textbf{\textit{#1}}}}
\newcommand{\ErrorTok}[1]{\textcolor[rgb]{0.64,0.00,0.00}{\textbf{#1}}}
\newcommand{\ExtensionTok}[1]{#1}
\newcommand{\FloatTok}[1]{\textcolor[rgb]{0.00,0.00,0.81}{#1}}
\newcommand{\FunctionTok}[1]{\textcolor[rgb]{0.13,0.29,0.53}{\textbf{#1}}}
\newcommand{\ImportTok}[1]{#1}
\newcommand{\InformationTok}[1]{\textcolor[rgb]{0.56,0.35,0.01}{\textbf{\textit{#1}}}}
\newcommand{\KeywordTok}[1]{\textcolor[rgb]{0.13,0.29,0.53}{\textbf{#1}}}
\newcommand{\NormalTok}[1]{#1}
\newcommand{\OperatorTok}[1]{\textcolor[rgb]{0.81,0.36,0.00}{\textbf{#1}}}
\newcommand{\OtherTok}[1]{\textcolor[rgb]{0.56,0.35,0.01}{#1}}
\newcommand{\PreprocessorTok}[1]{\textcolor[rgb]{0.56,0.35,0.01}{\textit{#1}}}
\newcommand{\RegionMarkerTok}[1]{#1}
\newcommand{\SpecialCharTok}[1]{\textcolor[rgb]{0.81,0.36,0.00}{\textbf{#1}}}
\newcommand{\SpecialStringTok}[1]{\textcolor[rgb]{0.31,0.60,0.02}{#1}}
\newcommand{\StringTok}[1]{\textcolor[rgb]{0.31,0.60,0.02}{#1}}
\newcommand{\VariableTok}[1]{\textcolor[rgb]{0.00,0.00,0.00}{#1}}
\newcommand{\VerbatimStringTok}[1]{\textcolor[rgb]{0.31,0.60,0.02}{#1}}
\newcommand{\WarningTok}[1]{\textcolor[rgb]{0.56,0.35,0.01}{\textbf{\textit{#1}}}}
\usepackage{graphicx}
\makeatletter
\newsavebox\pandoc@box
\newcommand*\pandocbounded[1]{% scales image to fit in text height/width
  \sbox\pandoc@box{#1}%
  \Gscale@div\@tempa{\textheight}{\dimexpr\ht\pandoc@box+\dp\pandoc@box\relax}%
  \Gscale@div\@tempb{\linewidth}{\wd\pandoc@box}%
  \ifdim\@tempb\p@<\@tempa\p@\let\@tempa\@tempb\fi% select the smaller of both
  \ifdim\@tempa\p@<\p@\scalebox{\@tempa}{\usebox\pandoc@box}%
  \else\usebox{\pandoc@box}%
  \fi%
}
% Set default figure placement to htbp
\def\fps@figure{htbp}
\makeatother
\setlength{\emergencystretch}{3em} % prevent overfull lines
\providecommand{\tightlist}{%
  \setlength{\itemsep}{0pt}\setlength{\parskip}{0pt}}
\usepackage{bookmark}
\IfFileExists{xurl.sty}{\usepackage{xurl}}{} % add URL line breaks if available
\urlstyle{same}
\hypersetup{
  pdftitle={Analyse en Composantes Principales (ACP)},
  pdfauthor={Projet ENSAssuRances},
  hidelinks,
  pdfcreator={LaTeX via pandoc}}

\title{Analyse en Composantes Principales (ACP)}
\usepackage{etoolbox}
\makeatletter
\providecommand{\subtitle}[1]{% add subtitle to \maketitle
  \apptocmd{\@title}{\par {\large #1 \par}}{}{}
}
\makeatother
\subtitle{Volet 4 - Apprentissage Non Supervisé}
\author{Projet ENSAssuRances}
\date{}

\begin{document}
\maketitle

{
\setcounter{tocdepth}{3}
\tableofcontents
}
\section{1. Introduction à l'Analyse
Multivariée}\label{introduction-uxe0-lanalyse-multivariuxe9e}

Dans cette étape, nous quittons la statistique univariée pour explorer
les relations complexes entre de multiples variables
\emph{simultanément}.

L'\textbf{Analyse en Composantes Principales (ACP)} est une méthode de
Machine Learning (apprentissage non supervisé) qui va nous permettre de
:

\begin{itemize}
\tightlist
\item
  \textbf{Résumer l'information} contenue dans notre grand nombre de
  variables quantitatives.
\item
  \textbf{Réduire la dimensionnalité} de notre base de données en créant
  de nouveaux axes de lecture (les composantes).
\item
  \textbf{Identifier des profils types} d'assurés et de véhicules en
  observant les corrélations directes (ex: un véhicule lourd est-il
  toujours associé à une prime d'assurance plus élevée ?).
\end{itemize}

\subsection{1.1. Préparation des données
actives}\label{pruxe9paration-des-donnuxe9es-actives}

L'algorithme de l'ACP obéit à des règles mathématiques strictes : il
nécessite des données \textbf{strictement numériques (continues)} et
\textbf{sans valeurs manquantes (NA)}.

Nous isolons donc les trois grandes familles de caractéristiques de
notre portefeuille :

\begin{itemize}
\tightlist
\item
  👤 \textbf{Le profil du conducteur} : Âge de l'assuré et ancienneté du
  permis de conduire.
\item
  🚗 \textbf{Le véhicule assuré} : Âge du véhicule, Poids, Puissance
  (DIN) et Valeur estimée.
\item
  📄 \textbf{Le contrat et le risque} : Montant de la franchise, Prime
  de base et Nombre d'antécédents de sinistres.
\end{itemize}

\begin{Shaded}
\begin{Highlighting}[]
\CommentTok{\# Chargement des librairies}
\FunctionTok{library}\NormalTok{(dplyr)}
\FunctionTok{library}\NormalTok{(FactoMineR)}
\FunctionTok{library}\NormalTok{(factoextra)}

\CommentTok{\# Chargement de la base globale}
\NormalTok{contrats\_globaux }\OtherTok{\textless{}{-}} \FunctionTok{readRDS}\NormalTok{(}\StringTok{"contrats\_exploration\_ok.rds"}\NormalTok{)}

\CommentTok{\# Sélection des variables continues pertinentes}
\NormalTok{data\_acp }\OtherTok{\textless{}{-}}\NormalTok{ contrats\_globaux }\SpecialCharTok{\%\textgreater{}\%}
  \FunctionTok{select}\NormalTok{(}
\NormalTok{    drv1Age, drv1DriveLicenceAge,    }\CommentTok{\# Profil du conducteur (Âge et Permis)}
\NormalTok{    vhAge, vhWeight, vhDIN, vhValue, }\CommentTok{\# Caractéristiques du véhicule (Âge, Poids, Puissance, Valeur)}
\NormalTok{    ctDeduc, COT\_AssBase, claimsAnt  }\CommentTok{\# Données du contrat (Franchise, Prime de base, Antécédents)}
\NormalTok{  ) }\SpecialCharTok{\%\textgreater{}\%}
\NormalTok{  tidyr}\SpecialCharTok{::}\FunctionTok{drop\_na}\NormalTok{() }\CommentTok{\# Suppression des lignes contenant des NA}

\CommentTok{\# Vérification}
\FunctionTok{print}\NormalTok{(}\StringTok{"{-}{-}{-} DIMENSIONS AVANT L\textquotesingle{}ACP {-}{-}{-}"}\NormalTok{)}
\end{Highlighting}
\end{Shaded}

\begin{verbatim}
## [1] "--- DIMENSIONS AVANT L'ACP ---"
\end{verbatim}

\begin{Shaded}
\begin{Highlighting}[]
\FunctionTok{cat}\NormalTok{(}\StringTok{"Nombre d\textquotesingle{}individus (lignes) :"}\NormalTok{, }\FunctionTok{nrow}\NormalTok{(data\_acp), }\StringTok{"}\SpecialCharTok{\textbackslash{}n}\StringTok{"}\NormalTok{)}
\end{Highlighting}
\end{Shaded}

\begin{verbatim}
## Nombre d'individus (lignes) : 301437
\end{verbatim}

\begin{Shaded}
\begin{Highlighting}[]
\FunctionTok{cat}\NormalTok{(}\StringTok{"Nombre de variables (colonnes) :"}\NormalTok{, }\FunctionTok{ncol}\NormalTok{(data\_acp), }\StringTok{"}\SpecialCharTok{\textbackslash{}n}\StringTok{"}\NormalTok{)}
\end{Highlighting}
\end{Shaded}

\begin{verbatim}
## Nombre de variables (colonnes) : 9
\end{verbatim}

\subsection{1.2. Exécution de l'algorithme et
Variances}\label{exuxe9cution-de-lalgorithme-et-variances}

Nous appliquons l'algorithme d'Analyse en Composantes Principales via la
fonction \texttt{PCA()} du package \texttt{FactoMineR}. Une étape
cruciale de ce calcul est le paramètre \texttt{scale.unit\ =\ TRUE}. En
effet, nos variables ayant des unités et des échelles radicalement
différentes (des années, des kilos, des euros), il est indispensable de
les centrer et de les réduire pour qu'elles aient toutes le même poids
dans le modèle.

Le premier résultat graphique à analyser est l'\textbf{éboulis des
valeurs propres}. Il nous indique le pourcentage d'information (la
variance) résumé par chaque nouvelle dimension créée.

\begin{Shaded}
\begin{Highlighting}[]
\CommentTok{\# Exécution de l\textquotesingle{}ACP }
\NormalTok{res\_acp }\OtherTok{\textless{}{-}} \FunctionTok{PCA}\NormalTok{(data\_acp, }\AttributeTok{scale.unit =} \ConstantTok{TRUE}\NormalTok{, }\AttributeTok{ncp =} \DecValTok{5}\NormalTok{, }\AttributeTok{graph =} \ConstantTok{FALSE}\NormalTok{)}

\CommentTok{\# Visualisation de la variance expliquée (Éboulis des valeurs propres)}
\NormalTok{graph\_vp }\OtherTok{\textless{}{-}} \FunctionTok{fviz\_eig}\NormalTok{(res\_acp, }
                     \AttributeTok{addlabels =} \ConstantTok{TRUE}\NormalTok{, }
                     \AttributeTok{ylim =} \FunctionTok{c}\NormalTok{(}\DecValTok{0}\NormalTok{, }\DecValTok{40}\NormalTok{),}
                     \AttributeTok{barfill =} \StringTok{"\#3498db"}\NormalTok{, }
                     \AttributeTok{barcolor =} \StringTok{"\#2980b9"}\NormalTok{,}
                     \AttributeTok{main =} \StringTok{"Éboulis des valeurs propres (Variance expliquée)"}\NormalTok{,}
                     \AttributeTok{xlab =} \StringTok{"Dimensions (Axes principaux)"}\NormalTok{,}
                     \AttributeTok{ylab =} \StringTok{"Pourcentage de variance"}\NormalTok{)}

\CommentTok{\# Affichage du graphique}
\FunctionTok{print}\NormalTok{(graph\_vp)}
\end{Highlighting}
\end{Shaded}

\begin{center}\includegraphics{04_ACP_files/figure-latex/execution-acp-1} \end{center}

\begin{Shaded}
\begin{Highlighting}[]
\FunctionTok{print}\NormalTok{(}\FunctionTok{head}\NormalTok{(res\_acp}\SpecialCharTok{$}\NormalTok{eig))}
\end{Highlighting}
\end{Shaded}

\begin{verbatim}
##        eigenvalue percentage of variance cumulative percentage of variance
## comp 1  2.0225421              22.472690                          22.47269
## comp 2  1.8055625              20.061806                          42.53450
## comp 3  1.2770541              14.189490                          56.72399
## comp 4  1.0008731              11.120812                          67.84480
## comp 5  0.9991371              11.101524                          78.94632
## comp 6  0.8618801               9.576445                          88.52277
\end{verbatim}

\subsection{1.3. Interprétation de la variance
expliquée}\label{interpruxe9tation-de-la-variance-expliquuxe9e}

Le tableau des valeurs propres nous permet de sélectionner le nombre
d'axes pertinents à retenir pour notre analyse :

\begin{itemize}
\tightlist
\item
  \textbf{Application du critère de Kaiser :} Les 4 premières
  composantes principales affichent une valeur propre (eigenvalue)
  supérieure à 1. Cela indique que l'information n'est pas concentrée
  sur un seul facteur prédominant, mais structurée autour de plusieurs
  dimensions.
\item
  \textbf{Qualité de représentation du premier plan :} L'Axe 1 (22,47
  \%) et l'Axe 2 (20,06 \%) capturent à eux seuls \textbf{42,53 \% de
  l'inertie totale} du nuage de points. Cette projection en deux
  dimensions (le premier plan factoriel) est de très bonne qualité pour
  des données comportementales et tarifaires issues du monde réel, et
  sera donc privilégiée pour lire nos corrélations.
\end{itemize}

\subsection{1.4. Le Cercle des Corrélations
(Variables)}\label{le-cercle-des-corruxe9lations-variables}

Le cercle des corrélations projette nos 9 variables initiales sur le
premier plan factoriel (Axes 1 et 2). * Des flèches proches (formant un
angle aigu) indiquent des variables fortement corrélées positivement. *
Des flèches opposées indiquent une corrélation négative. * Des flèches
formant un angle droit (90°) indiquent des variables indépendantes.

\begin{Shaded}
\begin{Highlighting}[]
\CommentTok{\# Création du cercle des corrélations avec factoextra}
\NormalTok{graph\_var }\OtherTok{\textless{}{-}} \FunctionTok{fviz\_pca\_var}\NormalTok{(res\_acp, }
             \AttributeTok{col.var =} \StringTok{"contrib"}\NormalTok{, }\CommentTok{\# On colore selon la contribution de la variable à l\textquotesingle{}axe}
             \AttributeTok{gradient.cols =} \FunctionTok{c}\NormalTok{(}\StringTok{"\#00AFBB"}\NormalTok{, }\StringTok{"\#E7B800"}\NormalTok{, }\StringTok{"\#FC4E07"}\NormalTok{), }\CommentTok{\# Du bleu (faible) au rouge (fort)}
             \AttributeTok{repel =} \ConstantTok{TRUE}\NormalTok{, }\CommentTok{\# Évite la superposition des étiquettes (très pratique !)}
             \AttributeTok{title =} \StringTok{"Cercle des Corrélations des variables d\textquotesingle{}Assurance (Axes 1 \& 2)"}\NormalTok{)}

\FunctionTok{print}\NormalTok{(graph\_var)}
\end{Highlighting}
\end{Shaded}

\begin{center}\includegraphics{04_ACP_files/figure-latex/cercle-correlations-1} \end{center}

\begin{Shaded}
\begin{Highlighting}[]
\CommentTok{\# ASTUCE POUR NOTRE BINÔME : On extrait les coordonnées exactes pour les analyser ensemble}
\FunctionTok{print}\NormalTok{(}\StringTok{"{-}{-}{-} COORDONNÉES DES VARIABLES SUR L\textquotesingle{}AXE 1 {-}{-}{-}"}\NormalTok{)}
\end{Highlighting}
\end{Shaded}

\begin{verbatim}
## [1] "--- COORDONNÉES DES VARIABLES SUR L'AXE 1 ---"
\end{verbatim}

\begin{Shaded}
\begin{Highlighting}[]
\FunctionTok{print}\NormalTok{(}\FunctionTok{round}\NormalTok{(res\_acp}\SpecialCharTok{$}\NormalTok{var}\SpecialCharTok{$}\NormalTok{coord[, }\DecValTok{1}\NormalTok{], }\DecValTok{2}\NormalTok{))}
\end{Highlighting}
\end{Shaded}

\begin{verbatim}
##             drv1Age drv1DriveLicenceAge               vhAge            vhWeight 
##                0.99                0.99                0.14               -0.02 
##               vhDIN             vhValue             ctDeduc         COT_AssBase 
##               -0.04               -0.08                0.00                0.21 
##           claimsAnt 
##               -0.01
\end{verbatim}

\begin{Shaded}
\begin{Highlighting}[]
\FunctionTok{print}\NormalTok{(}\StringTok{"{-}{-}{-} COORDONNÉES DES VARIABLES SUR L\textquotesingle{}AXE 2 {-}{-}{-}"}\NormalTok{)}
\end{Highlighting}
\end{Shaded}

\begin{verbatim}
## [1] "--- COORDONNÉES DES VARIABLES SUR L'AXE 2 ---"
\end{verbatim}

\begin{Shaded}
\begin{Highlighting}[]
\FunctionTok{print}\NormalTok{(}\FunctionTok{round}\NormalTok{(res\_acp}\SpecialCharTok{$}\NormalTok{var}\SpecialCharTok{$}\NormalTok{coord[, }\DecValTok{2}\NormalTok{], }\DecValTok{2}\NormalTok{))}
\end{Highlighting}
\end{Shaded}

\begin{verbatim}
##             drv1Age drv1DriveLicenceAge               vhAge            vhWeight 
##                0.06                0.06               -0.22                0.64 
##               vhDIN             vhValue             ctDeduc         COT_AssBase 
##                0.84                0.78                0.01                0.13 
##           claimsAnt 
##                0.00
\end{verbatim}

\subsection{1.5. Interprétation des Axes Factoriels (Conclusions
Métier)}\label{interpruxe9tation-des-axes-factoriels-conclusions-muxe9tier}

La projection des variables sur le premier plan factoriel révèle une
structure bipartite extrêmement nette dans le portefeuille
d'ENSAssuRances :

\begin{itemize}
\tightlist
\item
  \textbf{Axe 1 (L'Expérience du Conducteur) :} Cet axe horizontal est
  exclusivement défini par l'âge de l'assuré (coordonnée de 0.99) et
  l'ancienneté de son permis de conduire (0.99). Ces deux variables sont
  logiquement redondantes et parfaitement corrélées positivement.
\item
  \textbf{Axe 2 (La Gamme du Véhicule) :} Cet axe vertical capte les
  caractéristiques techniques et financières de la voiture. Il oppose
  les véhicules puissants (DIN : 0.84), chers (Value : 0.78) et lourds
  (Weight : 0.64) aux véhicules plus modestes et plus anciens (Age du
  véhicule : -0.22).
\item
  \textbf{Indépendance des profils :} Fait remarquable pour le
  positionnement commercial, ces deux groupes de variables sont
  orthogonaux (forment un angle de 90°). Dans cette base de données, la
  gamme du véhicule conduit est \textbf{indépendante} de l'âge du
  conducteur.
\item
  \textbf{Variables résiduelles :} Le montant de la franchise, les
  antécédents de sinistres et la prime de base sont très proches du
  centre de gravité (le point 0,0). Leurs variances ne sont pas
  capturées par ce premier plan en deux dimensions (Axes 1 et 2),
  suggérant que le risque et la tarification répondent à d'autres
  logiques multivariées ou catégorielles que nous explorerons avec
  l'ACM.
\end{itemize}

\subsection{1.5. Répartition des Individus et
Contributions}\label{ruxe9partition-des-individus-et-contributions}

\subsubsection{A. La Carte de Densité des Profils
(Individus)}\label{a.-la-carte-de-densituxe9-des-profils-individus}

Au lieu d'afficher 300 000 points superposés qui rendraient le graphique
illisible, nous utilisons une \textbf{carte de densité spatiale}. Plus
la couleur est foncée, plus la concentration d'assurés partageant le
même profil est forte.

\begin{Shaded}
\begin{Highlighting}[]
\FunctionTok{library}\NormalTok{(ggplot2)}
\FunctionTok{library}\NormalTok{(factoextra) }\CommentTok{\# On le recharge ici pour éviter l\textquotesingle{}erreur !}

\CommentTok{\# 1. Extraction des coordonnées exactes des 300 000 individus sur les axes 1 et 2}
\NormalTok{coord\_ind }\OtherTok{\textless{}{-}} \FunctionTok{as.data.frame}\NormalTok{(res\_acp}\SpecialCharTok{$}\NormalTok{ind}\SpecialCharTok{$}\NormalTok{coord)}

\CommentTok{\# 2. Création de la carte de densité (Heatmap 2D)}
\NormalTok{graph\_ind\_pro }\OtherTok{\textless{}{-}} \FunctionTok{ggplot}\NormalTok{(coord\_ind, }\FunctionTok{aes}\NormalTok{(}\AttributeTok{x =}\NormalTok{ Dim}\FloatTok{.1}\NormalTok{, }\AttributeTok{y =}\NormalTok{ Dim}\FloatTok{.2}\NormalTok{)) }\SpecialCharTok{+}
  \CommentTok{\# stat\_density\_2d crée de belles zones de couleurs selon la concentration}
  \FunctionTok{stat\_density\_2d}\NormalTok{(}\FunctionTok{aes}\NormalTok{(}\AttributeTok{fill =}\NormalTok{ ..level..), }\AttributeTok{geom =} \StringTok{"polygon"}\NormalTok{, }\AttributeTok{color =} \StringTok{"white"}\NormalTok{, }\AttributeTok{linewidth =} \FloatTok{0.1}\NormalTok{) }\SpecialCharTok{+}
  \FunctionTok{scale\_fill\_gradient}\NormalTok{(}\AttributeTok{low =} \StringTok{"\#ecf0f1"}\NormalTok{, }\AttributeTok{high =} \StringTok{"\#2c3e50"}\NormalTok{, }\AttributeTok{name =} \StringTok{"Densité}\SpecialCharTok{\textbackslash{}n}\StringTok{d\textquotesingle{}assurés"}\NormalTok{) }\SpecialCharTok{+}
  \FunctionTok{geom\_hline}\NormalTok{(}\AttributeTok{yintercept =} \DecValTok{0}\NormalTok{, }\AttributeTok{linetype =} \StringTok{"dashed"}\NormalTok{, }\AttributeTok{color =} \StringTok{"red"}\NormalTok{) }\SpecialCharTok{+}
  \FunctionTok{geom\_vline}\NormalTok{(}\AttributeTok{xintercept =} \DecValTok{0}\NormalTok{, }\AttributeTok{linetype =} \StringTok{"dashed"}\NormalTok{, }\AttributeTok{color =} \StringTok{"red"}\NormalTok{) }\SpecialCharTok{+}
  \FunctionTok{labs}\NormalTok{(}
    \AttributeTok{title =} \StringTok{"Concentration des profils d\textquotesingle{}assurés (Plan Factoriel)"}\NormalTok{,}
    \AttributeTok{subtitle =} \StringTok{"Lecture : Plus la zone est foncée, plus le nombre de clients est élevé"}\NormalTok{,}
    \AttributeTok{x =} \StringTok{"Axe 1 : Âge et Expérience du conducteur (22.5\%)"}\NormalTok{,}
    \AttributeTok{y =} \StringTok{"Axe 2 : Puissance et Valeur du véhicule (20.1\%)"}
\NormalTok{  ) }\SpecialCharTok{+}
  \FunctionTok{theme\_minimal}\NormalTok{()}

\FunctionTok{print}\NormalTok{(graph\_ind\_pro)}
\end{Highlighting}
\end{Shaded}

\begin{center}\includegraphics{04_ACP_files/figure-latex/graph-individus-densite-1} \end{center}

\subsubsection{B. Contribution des variables aux
axes}\label{b.-contribution-des-variables-aux-axes}

Ce graphique confirme mathématiquement notre lecture : quelles sont les
variables qui ``tirent'' le plus les individus vers la droite/gauche
(Axe 1) ou vers le haut/bas (Axe 2) ? La ligne rouge en pointillés
représente le seuil de contribution moyenne.

\begin{Shaded}
\begin{Highlighting}[]
\FunctionTok{library}\NormalTok{(gridExtra)}

\CommentTok{\# Barplot des contributions à l\textquotesingle{}Axe 1 (On utilise bien fviz\_contrib du package factoextra)}
\NormalTok{contrib\_axe1 }\OtherTok{\textless{}{-}} \FunctionTok{fviz\_contrib}\NormalTok{(res\_acp, }\AttributeTok{choice =} \StringTok{"var"}\NormalTok{, }\AttributeTok{axes =} \DecValTok{1}\NormalTok{, }
                             \AttributeTok{fill =} \StringTok{"\#3498db"}\NormalTok{, }\AttributeTok{color =} \StringTok{"black"}\NormalTok{, }
                             \AttributeTok{title =} \StringTok{"Ce qui définit l\textquotesingle{}Axe 1 (Âge)"}\NormalTok{)}

\CommentTok{\# Barplot des contributions à l\textquotesingle{}Axe 2}
\NormalTok{contrib\_axe2 }\OtherTok{\textless{}{-}} \FunctionTok{fviz\_contrib}\NormalTok{(res\_acp, }\AttributeTok{choice =} \StringTok{"var"}\NormalTok{, }\AttributeTok{axes =} \DecValTok{2}\NormalTok{, }
                             \AttributeTok{fill =} \StringTok{"\#e67e22"}\NormalTok{, }\AttributeTok{color =} \StringTok{"black"}\NormalTok{, }
                             \AttributeTok{title =} \StringTok{"Ce qui définit l\textquotesingle{}Axe 2 (Véhicule)"}\NormalTok{)}

\CommentTok{\# Affichage côte à côte}
\FunctionTok{grid.arrange}\NormalTok{(contrib\_axe1, contrib\_axe2, }\AttributeTok{ncol =} \DecValTok{2}\NormalTok{)}
\end{Highlighting}
\end{Shaded}

\begin{center}\includegraphics{04_ACP_files/figure-latex/unnamed-chunk-3-1} \end{center}

\end{document}
